% !TEX root = ../jlt.tex
\section{The compiler's own functional language}
\label{app:fjlt:ssa}

This appendix is the argument for the right-hand leg of Figure~\ref{fig:arrows}: that for a closed term $t$ of base type in the image of $\llbracket\cdot\rrbracket$, $\mathit{run}\,(\texttt{clang++}\,(\mathit{embed}\,t))$ is $\mathit{eval}\,t$.
Theorem~\ref{thm:simulation} gives that inside the calculus.
Here the program the compiler emits is read as the same term.
It is a reason, not a theorem: the compiler is not among the hypotheses, and its lowering is trusted, as every \Cpp{} program trusts it.

\paragraph{From the term to SSA.}
Instantiation is partial evaluation of the dictionary-passing term (Section~\ref{sec:binding}), and it leaves a closed first-order term over the base types.
The type abstractions are gone with their arguments, and the record parameters with the records the compiler resolved.
Every trait call is a call to one specialisation, and the only recursion left is $\mathit{iter}$ and $\mathit{fold}_F$.
For Go, Ellis et al.~\cite{ellis2022generic} prove the dictionary-passing translation correct beside monomorphisation, which Griesemer et al.~\cite{griesemer2020featherweight} proved; \Cpp{} only ever does the second, and the first carries the theorems.
The compiler lowers the first-order term to static single assignment form~\cite{cytron1991ssa}, and SSA is a functional language: Kelsey showed it to correspond to continuation-passing style~\cite{kelsey1995correspondence}, and Appel reads a basic block as a local function and its $\phi$-nodes as that function's parameters~\cite{appel1998ssa}.
So a Jlt function is translated twice into a functional language, by Appendix~\ref{app:fjlt:translation} into System~F with records and iterators and by the compiler into first-order functions with tail calls.
Rules~1 and~3 are what makes the source already look like the second.

\paragraph{A loop, read.}
The repository's IR fixture \texttt{showcase\_15} compiles Listing~\ref{lst:rule3}'s fold with its count \cppinline{n} known only at run time, at \texttt{-O2} with unrolling and vectorisation off.
Its one loop is
\begin{minted}{llvm}
6:                                    ; the loop header
  %7 = phi i64 [ %10, %6 ], [ %4, %3 ]    ; the count
  %8 = phi i64 [ %9, %6 ], [ 1, %3 ]      ; the vector,
  %9 = phi i64 [ %11, %6 ], [ 1, %3 ]     ;   two components
  %10 = add i64 %7, -1
  %11 = add i64 %9, %8
  %12 = icmp eq i64 %10, 0
  br i1 %12, label %13, label %6
\end{minted}
Its header's three $\phi$-nodes are the iterator's state in Figure~\ref{fig:fjlt-recursion}: the pair $(\llbracket r \rrbracket, \llbracket a \rrbracket)$, the chain's state its count and $\llbracket a \rrbracket$ the vector, itself a pair.
The body is the step, the matrix action with \cppinline{fib}'s entries as constants, reduced to one addition, $F(k+1) = F(k) + F(k-1)$, and the count's decrement ends the loop.
Nothing else survives: not the matrix, not the action's lambda, not the chain, not the vector as an object.
The function is the textbook loop for Fibonacci, which the source never wrote.
Both targets make the loop's state a parameter.
So Section~\ref{sec:binding}'s claim that the discipline costs nothing at run time has a structural ground: a Jlt fold names in the source the parameter the compiler would otherwise recover from assignments.
One seam shows here too.
SSA's values are the scalars in registers, and a value in memory, a boxed value in particular, is read and written through a pointer that no $\phi$-node carries.
That is the line between a value and a cell that Rule~1 draws.
